% Options for packages loaded elsewhere
\PassOptionsToPackage{unicode}{hyperref}
\PassOptionsToPackage{hyphens}{url}
\documentclass[
]{article}
\usepackage{xcolor}
\usepackage{amsmath,amssymb}
\setcounter{secnumdepth}{-\maxdimen} % remove section numbering
\usepackage{iftex}
\ifPDFTeX
  \usepackage[T1]{fontenc}
  \usepackage[utf8]{inputenc}
  \usepackage{textcomp} % provide euro and other symbols
\else % if luatex or xetex
  \usepackage{unicode-math} % this also loads fontspec
  \defaultfontfeatures{Scale=MatchLowercase}
  \defaultfontfeatures[\rmfamily]{Ligatures=TeX,Scale=1}
\fi
\usepackage{lmodern}
\ifPDFTeX\else
  % xetex/luatex font selection
\fi
% Use upquote if available, for straight quotes in verbatim environments
\IfFileExists{upquote.sty}{\usepackage{upquote}}{}
\IfFileExists{microtype.sty}{% use microtype if available
  \usepackage[]{microtype}
  \UseMicrotypeSet[protrusion]{basicmath} % disable protrusion for tt fonts
}{}
\makeatletter
\@ifundefined{KOMAClassName}{% if non-KOMA class
  \IfFileExists{parskip.sty}{%
    \usepackage{parskip}
  }{% else
    \setlength{\parindent}{0pt}
    \setlength{\parskip}{6pt plus 2pt minus 1pt}}
}{% if KOMA class
  \KOMAoptions{parskip=half}}
\makeatother
\usepackage{graphicx}
\makeatletter
\newsavebox\pandoc@box
\newcommand*\pandocbounded[1]{% scales image to fit in text height/width
  \sbox\pandoc@box{#1}%
  \Gscale@div\@tempa{\textheight}{\dimexpr\ht\pandoc@box+\dp\pandoc@box\relax}%
  \Gscale@div\@tempb{\linewidth}{\wd\pandoc@box}%
  \ifdim\@tempb\p@<\@tempa\p@\let\@tempa\@tempb\fi% select the smaller of both
  \ifdim\@tempa\p@<\p@\scalebox{\@tempa}{\usebox\pandoc@box}%
  \else\usebox{\pandoc@box}%
  \fi%
}
% Set default figure placement to htbp
\def\fps@figure{htbp}
\makeatother
\ifLuaTeX
  \usepackage{luacolor}
  \usepackage[soul]{lua-ul}
\else
  \usepackage{soul}
\fi
\setlength{\emergencystretch}{3em} % prevent overfull lines
\providecommand{\tightlist}{%
  \setlength{\itemsep}{0pt}\setlength{\parskip}{0pt}}
\usepackage{bookmark}
\IfFileExists{xurl.sty}{\usepackage{xurl}}{} % add URL line breaks if available
\urlstyle{same}
\hypersetup{
  pdftitle={Team \textless Team Name\textgreater{} - Mission \textless Mission Name\textgreater{}},
  hidelinks,
  pdfcreator={LaTeX via pandoc}}

\title{Team \hl{\textless Team Name\textgreater{}} - Mission
\hl{\textless Mission Name\textgreater{}}}
\author{}
\date{}

\begin{document}
\maketitle

\textbf{Abstract}

The Mission Report shall contain an Abstract. At a minimum, the abstract
shall identify the mission category in which the team is competing,
identify any unique/defining design characteristics of the project, and
provide whatever additional information may be necessary to convey any
other high-level project or program goals \& objectives.

\textbf{Authors}

First A. Author. \emph{Business or Academic Affiliation, State, Country,
Zip Code.}

Second B. Author. \emph{Business or Academic Affiliation, State,
Country, Zip Code.}

Third C. Author. \emph{Business or Academic Affiliation, State, Country,
Zip Code.}

Fourth D. Author. \emph{Business or Academic Affiliation, State,
Country, Zip Code.}

Nth N. Author. \emph{Business or Academic Affiliation, State, Country,
Zip Code}

\textbf{Nomenclature}

\emph{A} = amplitude of oscillation

\emph{a} = cylinder diameter

\emph{C\textsubscript{p}} = pressure coefficient

\emph{Cx} = force coefficient in the \emph{x} direction

\emph{Cy} = force coefficient in the \emph{y} direction

c = chord

d\emph{t} = time step

\begin{enumerate}
\def\labelenumi{\arabic{enumi}.}
\item
  \textbf{Introduction}
\end{enumerate}

The Mission Report shall contain an Introduction. This section provides
an overview of the academic program, stakeholders, team structure, and
mission organization and management strategies. The introduction may
repeat some of the content included in the abstract, because the
abstract is intended to act as a standalone synopsis if necessary.

\begin{enumerate}
\def\labelenumi{\arabic{enumi}.}
\setcounter{enumi}{1}
\item
  \textbf{Procedure for the Latin American Space Challenge Mission
  Report Submission}
\end{enumerate}

On or before the date specified in the Integrated Master Schedule prior
to the event, teams shall submit a single digital PDF copy of their
Mission Report. Mission reports shall be no longer than 50 pages. If
exceeding 10 Megabytes in size may need to be uploaded to a cloud server
if the permissions allow the judges unrestricted access to the document.

\begin{enumerate}
\def\labelenumi{\arabic{enumi}.}
\setcounter{enumi}{2}
\item
  \textbf{General Guidelines and Required Elements}
\end{enumerate}

The following section outlines general (non-formatting) guidelines to
follow if desired. Mission Reports will contain all elements identified
in the LASC LCSM, RCSM or SCSM documents, which is available for
download on the LASC website (https://www.lasc.space/documents). Missing
or incomplete elements will result in loss of points. The following
required items are provided for quick reference, and teams should
consult the LCSM, RCSM and SCSM for detailed guidance. Missions are
permitted to add other technical elements to their discretion. No table
of contents is required.

\subsection{Report Body}\label{report-body}

\begin{itemize}
\item
  Abstract
\item
  Introduction
\item
  System Architecture and Concept of Operations
\item
  Weights, Measures, and Performance Data
\item
  Computational Simulation
\item
  Conclusion and Lessons Learned
\end{itemize}

\subsection{Required Report
Appendices}\label{required-report-appendices}

\begin{itemize}
\item
  Hazard Analysis
\item
  Risk Assessment
\item
  Engineering Drawings
\item
  Optional Appendix
\end{itemize}

\subsection{List of References}\label{list-of-references}

\begin{itemize}
\item
  References are required and must follow the format under
  ``References'' at the end of this document.
\end{itemize}

\begin{enumerate}
\def\labelenumi{\arabic{enumi}.}
\setcounter{enumi}{3}
\item
  \textbf{Detailed Formatting Instructions}
\end{enumerate}

The styles and formats for the AIAA Papers Template have been
incorporated into the structure of this document. If you are using
Microsoft Word 2001 or later, please use this template to prepare your
manuscript. Regardless of which program you use to prepare your
manuscript, please use the formatting instructions contained in this
document as a guide.

\subsection{Document Text}\label{document-text}

The default font for the Project Technical Report is Times New Roman,
10-point size. In the electronic template, use the ``Text'' or
``Normal'' style from the pull-down menu to format all primary text for
your manuscript. The first line of every paragraph should be indented,
and all lines should be single-spaced. Default margins are 1'' on all
sides. In the electronic version of this template, all margins and other
formatting is preset. There should be no additional lines between
paragraphs.

\begin{quote}
Extended quotes, such as this example, are to be used when material
being cited is longer than a few sentences, or the standard quotation
format is not practical. In this Word template, the appropriate style is
``Extended Quote'' from the drop-down menu. Extended quotes are to be in
Times New Roman, 9-point font, indented 0.4'' and full justified.
\end{quote}

\emph{NOTE:} If you are using the electronic template to format your
manuscript, the required spacing and formatting will be applied
automatically, simply by using the appropriate style designation from
the pull-down menu.

\subsection{Headings}\label{headings}

The title of your paper should be typed in bold, 24-point type, with
capital and lower-case letters, and centered at the top of the page. The
names of the authors, business or academic affiliation, city, and
state/province should follow on separate lines below the title. The
names of authors with the same affiliation can be listed on the same
line above their collective affiliation information. Author names are
centered, and affiliations are centered and in italic type immediately
below the author names. The affiliation line for each author is to
include that author's city, state, and zip/postal code (or city,
province, zip/postal code and country, as appropriate). The first-page
footnotes (lower left-hand side) contain the job title and department
name, and street address/mail stop for each author. Author email
addresses may be included also.

Major headings (``Heading 1'' in the template style list) are bold
11-point font, flush-left, and numbered with numerals.

\begin{quote}
Subheadings (``Heading 2'' in the template style list) are bold, flush
left, and numbered with capital letters. Sub-Subheadings (``Heading 3''
in the template style list) are italic, flush left, and numbered (1. 2.
3. etc.)
\end{quote}

\subsection{Abstract}\label{abstract}

The abstract should appear at the beginning of your paper. It should be
one paragraph long (not an introduction) and complete (no reference
numbers). It should indicate subjects dealt with in the paper and state
the objectives of the investigation. Newly observed facts and
conclusions of the experiment or argument discussed in the paper must be
stated in summary form; readers should not have to read the paper to
understand the abstract. The abstract should be bold, indented 3 picas
(1/2'') on each side, and separated from the rest of the document by -
blank lines above and below the abstract text.

\subsection{Nomenclature}\label{nomenclature}

Papers with many symbols may benefit from a nomenclature list that
defines all symbols with units, inserted between the abstract and the
introduction. If one is used, it must contain all the symbology used in
the manuscript, and the definitions should not be repeated in the text.
In all cases, identify the symbols used if they are not widely
recognized in the profession. Define acronyms in the text, not in the
nomenclature.

\subsection{Footnotes and References}\label{footnotes-and-references}

Footnotes, where they appear, should be placed above the 1'' margin at
the bottom of the page. To insert footnotes into the template, use the
Insert\textgreater Footnote feature from the main menu as necessary.
Numbered footnotes as formatted automatically in the template are
acceptable, but superscript symbols are the preferred AIAA style, in the
sequence, *, †, ‡, §, ¶, \#, **. ††, ‡‡, §§, etc.

List and number all references at the end of the paper. Corresponding
bracketed numbers are used to cite references in the text {[}1{]},
unless the citation is an integral part of the sentence (e.g., ``It is
shown in Ref. {[}2{]} that\ldots'') or follows a mathematical
expression: ``A\textsuperscript{2} + B = C (Ref. {[}3{]}).'' For
multiple citations, separate reference numbers with commas {[}4, 5{]},
or use a dash to show a range {[}6-8{]}. Reference citations in the text
should be in numerical order.

In the reference list, give all authors' names; do not use ``et
al\emph{.}'' unless there are more than 10 authors. Papers that have not
been published should be cited as ``unpublished''; papers that have been
submitted or accepted for publication should be cited as ``submitted for
publication.'' Private communications and personal websites should
appear as footnotes rather than in the reference list.

References should be cited according to the standard publication
reference style (for examples, see the ``References'' section of this
template). Never edit titles in references to conform to AIAA style of
spellings, abbreviations, etc. Names and locations of publishers should
be listed; month and year should be included for reports and papers. For
papers published in translation journals, please give the English
citation first, followed by the original foreign language citation.

\subsection{Images, Figures, and
Tables}\label{images-figures-and-tables}

All artwork, captions, figures, graphs, and tables will be reproduced
exactly as submitted. Be sure to position any figures, tables, graphs,
or pictures as you want them printed.

Do not insert your tables and figures in text boxes. Figures should have
no background, borders, or outlines. In the electronic template, use the
``Figure'' style from the pull-down formatting menu to type caption
text. You may also insert the caption by going to the References menu
and choosing Insert Caption. Make sure the label is ``Fig.,'' and type
your caption text in the box provided. Captions are bold with a single
tab (no hyphen or other character) between the figure number and figure
description.

\includegraphics[width=3.45833in,height=2.63194in,alt={Image}]{./media/image1.jpg}

\textbf{Fig. 1. Magnetization as a function of applied fields.}

Place figure captions below all figures; place table titles above the
tables. If your figure has multiple parts, include the labels ``a),''
``b),'' etc. below and to the left of each part, above the figure
caption. Please verify that the figures and tables you mention in the
text exist. \emph{Please do not include captions as part of the figures,
and do not put captions in separate text boxes linked to the figures}.
When citing a figure in the text, use the abbreviation ``Fig.'' except
at the beginning of a sentence. Do not abbreviate ``Table.'' Number each
different type of illustration (i.e., figures, tables, images)
sequentially with relation to other illustrations of the same type.

Figure axis labels are often a source of confusion. Use words rather
than symbols. As in the example to the right, write the quantity
``Magnetization'' rather than just ``M.'' Do not enclose units in
parenthesis, but rather separate them from the preceding text by commas.
Do not label axes only with units. As in Fig. 1, for example, write
``Magnetization, kA/m'' not just ``kA/m.'' Do not label axes with a
ratio of quantities and units. For example, write ``Temperature, K,''
not ``Temperature/K.''

Multipliers can be especially confusing. Write ``Magnetization, kA/m''
or ``Magnetization, 10\textsuperscript{3} A/m.'' Do not write
``Magnetization (A/m) x 1000'' because the reader would not then know
whether the top axis label in Fig. 1 meant 16000 A/m or 0.016 A/m.
Figure labels must be legible, and all text within figures should be
uniform in style and size, no smaller than 8-point type.

\subsection{Equations, Numbers, Symbols, and
Abbreviations}\label{equations-numbers-symbols-and-abbreviations}

Equations are centered and numbered consecutively, with equation numbers
in parentheses flush right, as in Eq. (1). Insert a blank line above and
below the equation. First use the equation editor to create the
equation. If you are using Microsoft Word, use either the Microsoft
Equation Editor or the MathType add-on
(\href{http://www.mathtype.com}{\ul{http://www.mathtype.com}}) for
equations in your paper, use the function
(Insert\textgreater Object\textgreater Create New\textgreater Microsoft
Equation \emph{or} MathType Equation) to insert it into the document.
Please note that ``Float over text'' should \emph{not} be selected. To
insert the equation into the document:

\begin{enumerate}
\def\labelenumi{\arabic{enumi})}
\item
  Select the ``Equation'' style from the pull-down formatting menu and
  hit ``tab'' once.
\item
  Insert the equation, hit ``tab'' again,
\item
  Enter the equation number in parentheses.
\end{enumerate}

A sample equation is included here, formatted using the preceding
instructions. To make your equation more compact, you can use the
solidus (/), the exp function, or appropriate exponents. Use parentheses
to avoid ambiguities in denominators.

\includegraphics[width=3.09028in,height=0.54861in,alt={Image2}]{./media/image2.jpg}
(1)

Be sure that the symbols in your equation are defined before the
equation appears or immediately following. Italicize symbols (\emph{T}
might refer to temperature, but T is the unit tesla). Refer to ``Eq.
(1),'' not ``(1)'' or ``equation (1)'' except at the beginning of a
sentence: ``Equation (1) is\ldots'' Equations can be labeled other than
``Eq.'' should they represent inequalities, matrices, or boundary
conditions. If what is represented is more than one equation, the
abbreviation ``Eqs.'' can be used.

Define abbreviations and acronyms the first time they are used in the
text, even after they have already been defined in the abstract. Very
common abbreviations such as AIAA, SI, ac, and dc do not have to be
defined. Abbreviations that incorporate periods should not have spaces:
write ``P.R.,'' not ``P. R.'' Delete periods between initials if the
abbreviation has three or more initials; e.g., U.N. but ESA. Do not use
abbreviations in the title unless they are unavoidable (for instance,
``AIAA'' in the title of this article).

\subsection{General Grammar and Preferred
Usage}\label{general-grammar-and-preferred-usage}

Use only one space after periods or colons. Hyphenate complex modifiers:
``zero-field-cooled magnetization.'' Avoid dangling participles, such
as, ``Using Eq. (1), the potential was calculated.'' {[}It is not clear
who or what used Eq. (1).{]} Write instead ``The potential was
calculated using Eq. (1),'' or ``Using Eq. (1), we calculated the
potential.''

Use a zero before decimal points: ``0.25,'' not ``.25.'' Use
``cm\textsuperscript{2},'' not ``cc.'' Indicate sample dimensions as
``0.1 cm x 0.2 cm,'' not ``0.1 x 0.2 cm\textsuperscript{2}.'' The
preferred abbreviation for ``seconds'' is ``s,'' not ``sec.'' Do not mix
complete spellings and abbreviations of units: use
``Wb/m\textsuperscript{2}'' or ``webers per square meter,'' not
``webers/m\textsuperscript{2}.'' When expressing a range of values,
write ``7 to 9'' or ``7-9,'' not ``7\textasciitilde9.''

A parenthetical statement at the end of a sentence is punctuated outside
of the closing parenthesis (like this). (A parenthetical sentence is
punctuated within parenthesis.) In American English, periods and commas
are placed within quotation marks, like ``this period.'' Other
punctuation is ``outside''! Avoid contractions; for example, write ``do
not'' instead of ``don't.'' The serial comma is preferred: ``A, B, and
C'' instead of ``A, B and C.''

If you wish, you may write in the first person singular or plural and
use the active voice (``I observed that\ldots'' or ``We observed
that\ldots'' instead of ``It was observed that\ldots''). Remember to
check spelling. If your native language is not English, please ask a
native English-speaking colleague to proofread your paper.

The word ``data'' is plural, not singular (i.e., ``data are,'' not
``data is''). The subscript for the permeability of vacuum
µ\textsubscript{0} is zero, not a lowercase letter ``o.'' The term for
residual magnetization is ``remanence''; the adjective is ``remanent'';
do not write ``remnance'' or ``remnant.'' The word ``micrometer'' is
preferred over ``micron'' when spelling out this unit of measure. A
graph within a graph is an ``inset,'' not an ``insert.'' The word
``alternatively'' is preferred to the word ``alternately'' (unless you
really mean something that alternates). Use the word ``whereas'' instead
of ``while'' (unless you are referring to simultaneous events). Do not
use the word ``essentially'' to mean ``approximately'' or
``effectively.'' Do not use the word ``issue'' as a euphemism for
``problem.'' When compositions are not specified, separate chemical
symbols by en-dashes; for example, ``NiMn'' indicates the intermetallic
compound Ni\textsubscript{0.5}Mn\textsubscript{0.5} whereas ``Ni--Mn''
indicates an alloy of some composition
Ni\textsubscript{x}Mn\textsubscript{1-x}.

Be aware of the different meanings of the homophones ``affect'' (usually
a verb) and ``effect'' (usually a noun), ``complement'' and
``compliment,'' ``discreet'' and ``discrete,'' ``principal'' (e.g.,
``principal investigator'') and ``principle'' (e.g., ``principle of
measurement''). Do not confuse ``imply'' and ``infer.''

Prefixes such as ``non,'' ``sub,'' ``micro,'' ``multi,'' and ``"ultra''
are not independent words; they should be joined to the words they
modify, usually without a hyphen. There is no period after the ``et'' in
the abbreviation ``et al\emph{.}'' The abbreviation ``i.e.,'' means
``that is,'' and the abbreviation ``e.g.,'' means ``for example'' (these
abbreviations are not italicized).

\begin{enumerate}
\def\labelenumi{\arabic{enumi}.}
\setcounter{enumi}{4}
\item
  \textbf{Conclusions and Lessons Learned}
\end{enumerate}

The Mission Report shall contain Conclusions and Lessons Learned. This
section shall include the lessons learned during the design,
manufacturing, and testing of the project, both from a team management
and technical development perspective. Furthermore, this section should
include strategies for corporate knowledge transfer from senior team
members to the rising underclassmen who will soon take their place.

\begin{enumerate}
\def\labelenumi{\arabic{enumi}.}
\setcounter{enumi}{5}
\item
  \textbf{References}
\end{enumerate}

The following pages are intended to provide examples of the different
reference types. All references should be in 9-point font, with the
first line flush left and reference numbers inserted in brackets. You
are not required to indicate the type of reference; different types are
shown here for illustrative purposes only. The DOI (digital object
identifier) should be incorporated in every reference for which it is
available (see Ref. 1 sample); for more information on DOIs, visit
www.doi.org or \href{http://www.crossref.org}{www.crossref.org}.

\subsection{Periodicals}\label{periodicals}

{[}1{]} Vatistas, G. H., Lin, S., and Kwok, C. K., ``Reverse Flow Radius
in Vortex Chambers,'' \emph{AIAA Journal}, Vol. 24, No. 11, 1986, pp.
1872, 1873.

doi: 10.2514/3.13046

{[}2{]} Alyanak, E. J., and Pendleton, E., ``Aeroelastic Tailoring and
Active Aeroelastic Wing Impact on a Lambda Wing Configuration,'' Journal
of Aircraft, published online 10 Nov. 2016.

doi: 10.2514/1.C033040

{[}3{]} Dornheim, M. A., ``Planetary Flight Surge Faces Budget
Realities,'' \emph{Aviation Week and Space Technology}, Vol. 145, No.
24, 9 Dec. 1996, pp. 44--46.

{[}4{]} Terster, W., ``NASA Considers Switch to Delta 2,'' \emph{Space
News}, Vol. 8, No. 2, 13--19 Jan. 1997, pp. 1, 18.

All of the preceding information is required. The journal issue number
(``No. 11'' in Ref. 1) is preferred, but the month (Nov.) can be
substituted if the issue number is not available. Use the complete date
for daily and weekly publications. Transactions follow the same style as
other journals.

\subsection{Books}\label{books}

{[}5{]} Peyret, R., and Taylor, T. D., \emph{Computational Methods in
Fluid Flow}, 2\textsuperscript{nd} ed., Springer-Verlag, New York, 1983,
Chaps. 7, 14.

{[}6{]} Oates, G. C. (ed.), \emph{Aerothermodynamics of Gas Turbine and
Rocket Propulsion}, AIAA Education Series, AIAA, New York, 1984, pp. 19,
136.

{[}7{]} Volpe, R., ``Techniques for Collision Prevention, Impact
Stability, and Force Control by Space Manipulators,''
\emph{Teleoperation and Robotics in Space}, edited by S. B. Skaar and C.
F. Ruoff, Progress in Astronautics and Aeronautics, AIAA, Washington,
DC, 1994, pp. 175--212.

Publisher, place, and date of publication are required for all books. No
state or country is required for major cities: New York, London, Moscow,
etc. A differentiation must always be made between Cambridge, MA, and
Cambridge, England, UK. Note that series titles are in Roman type.

\subsection{Proceedings}\label{proceedings}

{[}8{]} Thompson, C. M., ``Spacecraft Thermal Control, Design, and
Operation,'' \emph{AIAA Guidance, Navigation, and Control Conference},
CP849, Vol. 1, AIAA, Washington, DC, 1989, pp. 103--115

{[}9{]} Chi, Y. (ed.), \emph{Fluid Mechanics Proceedings}, NASA SP-255,
1993.

{[}10{]} Morris, J. D., ``Convective Heat Transfer in Radially Rotating
Ducts,'' \emph{Proceedings of the Annual Heat Transfer Conference},
edited by B. Corbell, Vol. 1, Inst. of Mechanical Engineering, New York,
1992, pp. 227--234.

\subsection{Reports, Theses, and Individual
Papers}\label{reports-theses-and-individual-papers}

{[}11{]} Chapman, G. T., and Tobak, M., ``Nonlinear Problems in Flight
Dynamics,'' NASA TM-85940, 1984.

{[}12{]} Brandis, A. M., Johnston, C. O., and Cruden, B. A.,
``Nonequilibrium Radiation for Earth Entry,'' AIAA Paper 2016-3690, June
2016.

{[}13{]} Steger, J. L., Jr., Nietubicz, C. J., and Heavey, J. E., ``A
General Curvilinear Grid Generation Program for Projectile
Configurations,'' U.S. Army Ballistic Research Lab., Rept.
ARBRL-MR03142, Aberdeen Proving Ground, MD, Oct. 1981.

{[}14{]} Tseng, K., ``Nonlinear Green's Function Method for Transonic
Potential Flow,'' Ph.D. Dissertation, Aeronautics and Astronautics
Dept., Boston Univ., Cambridge, MA, 1983.

Government agency reports do not require locations. For reports such as
NASA TM-85940, neither insert nor delete dashes; leave them as provided.
Place of publication \emph{should} be given, although it is not
mandatory, for military and company reports. Always include a city and
state for universities. Papers need only the name of the sponsor;
neither the sponsor's location nor the conference name and location is
required. \emph{Do not confuse proceedings references with conference
papers}.

\subsection{Electronic Publications}\label{electronic-publications}

Regularly issued electronic journals and other publications are
permitted as references. Include the DOI if provided; otherwise provide
the full URL. Archived data sets also may be referenced if the material
is openly accessible, and the repository is committed to achieving the
data indefinitely. References to electronic data available only from
personal websites or commercial, academic, or government ones where
there is no commitment to archiving the data are not permitted in the
reference list.

{[}15{]} Atkins, C. P., and Scantelbury, J. D., ``The Activity
Coefficient of Sodium Chloride in a Simulated Pore Solution
Environment,'' \emph{Journal of Corrosion Science and Engineering}
{[}online journal{]}, Vol. 1, No. 1, Paper 2, URL:
\href{http://www.cp/umist.ac.uk/JCSE/vol1/vol1.html}{\ul{http://www.cp/umist.ac.uk/JCSE/vol1/vol1.html}}
{[}retrieved 13 April 1998{]}.

{[}16{]} Vickers, A., ``10-110 mm/hr Hypodermic Gravity Design A,''
\emph{Rainfall Simulation Database} {[}online database{]}, URL:
\href{http://www.geog.le.ac.uk/bgrg/lab.htm}{\ul{http://www.geog.le.ac.uk/bgrg/lab.htm}}
{[}retrieved 15 March 2006{]}.

Break website addresses after punctuation, and do not hyphenate at line
breaks.

\emph{Computer Software}

{[}17{]} TAPP, Thermochemical and Physical Properties, Software Package,
Ver. 1.0, E. S. Microware, Hamilton, OH, 1992.

Include a version number and the company name and location of software
packages.

\emph{Patents}

Patents appear infrequently. Be sure to include the patent number and
date.

{[}18{]} Scherrer, R., Overholster, D., and Watson, K., Lockheed Corp.,
Burbank, CA, U.S. Patent Application for a ``Vehicle,'' Docket No.
P-01-1532, filed 11 Feb. 1979.

\emph{Private Communications and Websites}

References to private communications and personal website addresses are
not permitted. They may, however, be incorporated into the main text of
a manuscript or may appear in footnotes.

\emph{Unpublished Papers and Books}

Unpublished works can be used as references as long as they are being
considered for publication or can be located by the reader (such as
papers that are part of an archival collection). If a journal paper or a
book is being considered for publication, choose the format that
reflects the status of the work (depending upon whether it has been
accepted for publication):

{[}19{]} Doe, J., ``Title of Paper,'' \emph{Name of Journal} (to be
published).

{[}20{]} Doe, J., ``Title of Chapter,'' \emph{Name of Book}, edited
by\ldots, Publisher's name and location (to be published).

{[}21{]} Doe, J., ``Title of Work,'' Name of Archive, Univ. (or
organization), City, State, Year (unpublished).

Unpublished works in an archive \emph{must} include the name of the
archive and the name and location of the university or other
organization where the archive is held. Also include any cataloging
information that may be provided.

\textbf{\hfill\break
}

\textbf{Hazard Analysis Appendix}

The first Mission Report appendix shall contain a Hazard Analysis
Report. This appendix shall address as applicable, hazardous material
handling, transportation, storage procedures, and any other aspects of
the design which pose potential hazards to operating personnel. A
mitigation approach -- by process and/or design -- shall be defined for
each hazard identified.

\textbf{\hfill\break
}

\textbf{Risk Assessment Appendix}

The second Mission Report appendix shall contain a Risk Assessment. This
appendix shall summarize risk and reliability concepts associated with
the project. All identified failure modes which pose a risk to mission
success shall be recorded in a matrix, organized according to the
mission phases identified by CONOPS.

\textbf{\hfill\break
}

\textbf{Engineering Drawings and Optional Appendix}

The final Mission Report appendix shall contain Engineering Drawings and
optional appendices. This appendix shall include any revision controlled
technical drawings necessary to define significant subsystems or
components, and other optional appendices can include, but are not
limited to further Subsystem Details, Launch Support Equipment Details,
Detailed Structural and Mechanical Calculation, Detailed Logical Process
Diagrams, Detailed Software Architecture, and Detailed Electrical
Architecture.

\end{document}


